\section{The discriminant pairing and admissible columns}\label{sec:discriminant}

We prove Proposition~\ref{prop:admissible-potential}. The arithmetic
part of the proof colors isotropic elements of the discriminant group
so that the pairing on differences within one color has controlled
image order. The geometric part converts this order bound into a
dimension bound for a robust subset of admissible columns.

\subsection{Orthogonal blocks and a coloring of isotropic elements}

For a nonsingular symmetric integer matrix \(B\), the formula
\[
 \lambda_B(\bar x,\bar y)=x^{\mathsf T}B^{-1}y\bmod\Z
 \qquad(\bar x,\bar y\in G_B)
\]
defines a symmetric pairing with values in \(\mathbb Q/\Z\).
Changing \(x\) or \(y\) by a vector in \(B\Z^m\) changes the displayed
value by an integer. Moreover, if \(\lambda_B(\bar x,\bar y)=0\) for
every \(y\in\Z^m\), then \(B^{-1}x\in\Z^m\), so \(\bar x=0\).
Thus the induced map from \(G_B\) to
\(\Hom(G_B,\mathbb Q/\Z)\) is an isomorphism. We call such a pairing
\emph{perfect}.

Distinct primary components are orthogonal: a pairing value between
an element of \(p\)-power order and one of \(q\)-power order is
annihilated by coprime powers when \(p\ne q\). The restriction to each
primary component is consequently perfect. An element is
\emph{isotropic} if its self-pairing is zero. If an element of \(G_B\)
is isotropic, then so is its projection to each primary component,
by the uniqueness of primary decomposition in \(\mathbb Q/\Z\).
In particular the class of every \(z\in\mathcal S(B)\) is isotropic.

Orthogonal decompositions of these pairings are standard in the theory
of finite linking forms; see Miranda~\cite[Section~1]{Miranda}.
We include the elementary splitting argument and the precise block
forms needed here.

\begin{lemma}[Orthogonal block decomposition]\label{lem:discriminant-blocks}
Let \(G\) be a finite abelian \(p\)-group with a perfect symmetric
pairing \(\lambda:G\times G\to\mathbb Q/\Z\). It is an orthogonal
direct sum of blocks of the following forms:
\begin{enumerate}
 \item a cyclic group of order \(p^j\) whose generator has
 self-pairing \(a/p^j\bmod\Z\), with \(p\nmid a\);
 \item when \(p=2\), a group \((\Z/2^j\Z)^2\) with pairing matrix
 \[
  2^{-j}\begin{pmatrix}2a&c\\c&2d\end{pmatrix}\bmod\Z,
  \qquad c\ \text{odd}.
 \]
\end{enumerate}
\end{lemma}

\begin{proof}
Let \(p^j\) be the exponent of a nonzero group still to be split.
If some element has self-pairing of order \(p^j\), it has order
\(p^j\) and its cyclic subgroup is a block of the first kind.
Otherwise choose an element \(x\) of order \(p^j\). Perfectness
implies that the character \(\lambda(x,\cdot)\) has order \(p^j\),
so there is an element \(y\) for which \(\lambda(x,y)\) has order
\(p^j\). Write the three pairing values on \(x,y\), with denominator
\(p^j\), using numerators \(a,c,d\), respectively. Then \(c\) is
a unit modulo \(p\), while \(a,d\) are divisible by \(p\).
For odd \(p\), the self-pairing numerator of \(x+y\) is
\(a+2c+d\), a unit, so again a block of the first kind is available.

For \(p=2\), the two diagonal numerators are even and the
off-diagonal numerator is odd. The resulting two by two numerator
matrix is invertible modulo \(2^j\). Pairing a relation between
\(x,y\) with \(x,y\) therefore shows that its two coefficients
vanish modulo \(2^j\). Thus they span \((\Z/2^j\Z)^2\), and the
restriction of the pairing to this subgroup is perfect, giving the
second kind of block.

It remains to justify splitting a block \(H\) off orthogonally.
Because the pairing on \(H\) is perfect, for any \(g\in G\) there
is a unique \(h\in H\) such that
\(\lambda(g,\cdot)|_H=\lambda(h,\cdot)|_H\). Hence
\(g-h\in H^\perp\). Also \(H\cap H^\perp=0\), so
\(G=H\oplus H^\perp\). The pairing on \(H^\perp\) is perfect,
since an element annihilating both summands annihilates \(G\).
Iteration proves the decomposition.
\end{proof}

The exponents of the cyclic factors in this decomposition are the
invariant exponents \(b_1\ge b_2\ge\cdots\) of \(G\). A block of the
second kind contributes two equal exponents.

\begin{lemma}[Primary coloring]\label{lem:primary-coloring}
In the setting of Lemma~\ref{lem:discriminant-blocks}, put \(h=b_3\),
with zero exponents appended as necessary. In an orthogonal block
decomposition, let \(U\) be the sum of the blocks of exponent
strictly greater than \(h\), and let \(T\) be the remaining sum.
There is a coloring of the isotropic elements of \(G\) with \(C_p\)
colors such that elements \(x,y\) of one color satisfy
\begin{equation}\label{eq:primary-color-pairing}
 p^h\lambda_U(x_U,y_U)=0\quad\text{in }\mathbb Q/\Z.
\end{equation}
Here \(x_U,y_U\) denote the projections to \(U\). The color counts
may be chosen to satisfy
\[
 C_p\le 1+8b_1,\qquad
 C_p\le 2^{b_2-b_3}\quad\text{if \(p\) is odd}.
\]
\end{lemma}

The cutoff at \(b_3\) leaves at most two cyclic factors above it.
Up to a lower-order error, the color count will be charged to height-two
layers, while layers of height at least three will control the order of
the pairing image on differences within one color.

\begin{proof}
The exponent of \(T\) divides \(p^h\). Isotropy and orthogonality
therefore imply
\begin{equation}\label{eq:top-self-denominator}
 p^h\lambda_U(x_U,x_U)=0
\end{equation}
for every element under consideration. There are at most two cyclic
factors in \(U\). A block of the second kind cannot be divided by
the cutoff at \(h\), since its two factors have equal exponent.
We can thus handle all possibilities explicitly. In the following,
\(v_p(0)=+\infty\); for a nonzero cyclic coordinate its valuation is
that of any representative not divisible by the full modulus.

If \(U=0\), use one color. If it has one factor, this is a cyclic
block of exponent \(j>h\), with coordinate \(x\) and a unit
self-pairing numerator. Equation~\eqref{eq:top-self-denominator}
forces
\[
 v_p(x)\ge \left\lceil\frac{j-h}{2}\right\rceil.
\]
Two such coordinates have mutual pairing annihilated by \(p^h\),
so again one color suffices.

Suppose next that \(U\) consists of two orthogonal cyclic blocks of
exponents \(j\ge i>h\), with respective unit self-pairing numerators
\(a_1,a_2\). For their coordinates \(x,y\), put
\[
 F=\max\{j-2v_p(x),\,i-2v_p(y)\}.
\]
Use one \emph{deep} color for \(F\le h\). Two deep elements have the
required mutual denominator term by term, since the sum of the two
valuations in either coordinate is at least its exponent minus \(h\).
If \(F>h\), the two entries in the maximum must equal \(F\).
Otherwise the term of largest denominator in the self-pairing
could not cancel, contrary to
Equation~\eqref{eq:top-self-denominator}. In particular \(F\le i\),
and \(F\) has the parity of both \(i\) and \(j\). Write
\[
 x=p^{(j-F)/2}X,\qquad y=p^{(i-F)/2}Y,
\]
where \(X,Y\) are units. With \(n=F-h>0\), the self-pairing
condition is
\[
 a_1X^2+a_2Y^2\equiv0\pmod{p^n}.
\]
The residues of \(X,Y\) modulo \(p^n\) are well defined: they are
determined modulo \(p^{(j+F)/2}\) and \(p^{(i+F)/2}\), respectively,
and both exponents are at least \(F\ge n\).
Color this element by \(F\) and
\(r=XY^{-1}\bmod p^n\). Its ratio satisfies
\(r^2\equiv-a_2a_1^{-1}\pmod{p^n}\). For two elements with the
same color, using primes for the coordinates of the second one,
\[
 a_1XX'+a_2YY'
 \equiv YY'(a_1r^2+a_2)\equiv0\pmod{p^n}.
\]
This is precisely Equation~\eqref{eq:primary-color-pairing}.

For completeness, a unit square has at most two square roots modulo
an odd prime power and at most four modulo a power of \(2\).
Indeed, for two unit roots \(r,s\), the modulus divides
\((r-s)(r+s)\). At an odd prime at most one factor is divisible by
\(p\), giving \(r\equiv s\) or \(r\equiv-s\pmod{p^n}\). At \(2\),
for odd \(r,s\), one of \(r-s,r+s\) has valuation exactly one and
the other is divisible by four, with the convention of infinite
valuation at zero. For \(n\ge2\) this gives
\(r\equiv s\) or \(r\equiv-s\pmod{2^{n-1}}\), allowing at most four
residues modulo \(2^n\); \(n=1\) is immediate.

There are at most \(\lceil(i-h)/2\rceil\) possible nondeep values of
\(F\). At odd \(p\), the count is thus at most
\(1+2\lceil(i-h)/2\rceil\). For \(i-h\ge2\) this is at most
\(2^{i-h}\). If \(i-h=1\) and a nondeep element exists, its value
is \(F=h+1\). The parity just established shows that in a deep
element the two exponent-minus-twice-valuation quantities are at
most \(h-1\). In the coordinate of exponent \(j\), for example,
a deep valuation is at least \((j-h+1)/2\), whereas a nondeep
valuation is \((j-h-1)/2\). Their sum is at least \(j-h\), and
the same argument applies to the coordinate of exponent \(i\).
Every deep element therefore has the required mutual denominator
with every nondeep element. Merge the deep color into one nondeep
color, leaving at most two colors. If there is no nondeep element,
one color suffices. This proves the odd-prime bound in this case.
At \(2\), the unmerged count is at most
\(1+4\lceil(i-h)/2\rceil\), which satisfies the asserted
polynomial bound.

The remaining possibility is a single block of the second kind at
\(2\), of exponent \(j>h\). Write its numerator matrix as
\(\left(\begin{smallmatrix}2a&c\\c&2d\end{smallmatrix}\right)\)
with \(c\) odd, and put
\[
 t=\min\{v_2(x),v_2(y)\},\qquad F=j-2t.
\]
Again use a deep color for \(F\le h\). For two deep elements, the
valuation of each coordinate in each element is at least
\(\lceil(j-h)/2\rceil\), which proves the mutual denominator bound.
For \(F>h\), write \(x=2^tX,y=2^tY\), so at least one of \(X,Y\)
is odd, and put \(n=F-h\). Choose the first coordinate as pivot if
it is odd, and otherwise choose the second. After ordering the
pivot first, call it \(u\), call the other coordinate \(v\), and
write
\[
 f(X,Y)=aX^2+cXY+dY^2,
\]
interchanging \(a,d\) when the second coordinate is the pivot.
The self-pairing numerator is twice this quadratic form. Thus
Equation~\eqref{eq:top-self-denominator} says
\[
 f(u,v)\equiv0\pmod{2^{n-1}},
\]
where the condition is empty for \(n=1\). Color the element by
\(F\), the pivot, and \(r=vu^{-1}\bmod 2^n\). This ratio satisfies
\[
 f(1,r)=a+cr+dr^2\equiv0\pmod{2^{n-1}}.
\]
The ratio is well defined modulo \(2^n\), since the normalized
coordinates are determined modulo \(2^{j-t}\) and \(n\le j-t\).
Its derivative \(c+2dr\) is odd. Every root modulo \(2^\alpha\), for
\(\alpha\ge1\), therefore has exactly one root lift modulo
\(2^{\alpha+1}\): the two possible lifts change the value modulo
\(2^{\alpha+1}\) by \(2^\alpha\) times an odd number. There are at most two roots modulo
\(2\). It follows that there are at most four permissible ratios
modulo \(2^n\), including the one extra bit beyond the modulus in
the self-pairing condition. This bound also holds for \(n=1\).

For two elements of the same color, their normalized pivot
coordinates are odd numbers \(u,u'\) and their ratios satisfy
\(r'\equiv r\pmod{2^n}\). Their bilinear numerator modulo \(2^n\)
is
\[
 uu'\bigl(2a+c(r+r')+2drr'\bigr)
 \equiv 2uu'f(1,r)\equiv0\pmod{2^n}.
\]
Restoring the common factor \(2^{2t}\) proves
Equation~\eqref{eq:primary-color-pairing}. There are at most \(b_1\)
possible values of \(F>h\), two pivots, and four ratios at each
pivot, so the count is at most \(1+8b_1\). The same bound holds
in all preceding cases. Finally, the number of factors above \(h\)
is zero, one, or two according as handled above, and their second
exponent satisfies \(i-h=b_2-b_3\) when there are two. This proves
the asserted bounds.
\end{proof}

Apply Lemma~\ref{lem:primary-coloring} independently to the primary
components of \(G_B\). An overall color is the tuple of its primary
colors; let \(C_{\mathrm{col}}(B)\) be the product of the counts over
the nontrivial primary components.
The following explicit estimate is the color cost that we need:
\begin{equation}\label{eq:discriminant-color-cost}
 \frac{\log_2 C_{\mathrm{col}}(B)}{N}
 \le \frac54\sum_{\substack{p,e\\k_B(p,e)=2}}\nu_p
       +O(N^{-1/5}\log N)
 \le \beta\sum_{\substack{p,e\\k_B(p,e)=2}}\nu_p
       +O(N^{-1/5}\log N).
\end{equation}
Indeed, for \(p\ge N^{4/5}\) the prime is odd for large \(N\), and
the logarithm of its color count is at most \(b_{p,2}-b_{p,3}\).
This difference is exactly the number of its height-two layers,
while \(\nu_p\ge4/(5N)\). These primes give the first term of the
first bound. For \(p<N^{4/5}\), use instead
\(C_p\le1+8b_{p,1}\). The order bound for \(G_B\) gives
\[
 b_{p,1}\le\log_2|G_B|\le (N/2)\log_2N.
\]
There are at most \(N^{4/5}\) such primes, so they cost
\(O(N^{4/5}\log N)\) bits in total. This proves the first bound;
the second uses \(5/4<\beta=7/5\). The error term in
Equation~\eqref{eq:discriminant-color-cost} is
\(o(\ell^{-A})\) for every fixed \(A>0\).

Set
\begin{equation}\label{eq:discriminant-high-weight}
 J_B=\sum_{\substack{p,e\\k_B(p,e)\ge3}}k_B(p,e)\nu_p.
\end{equation}
Equation~\eqref{eq:group-order} gives \(0\le J_B\le1/2\).

\begin{lemma}[Pairing image on differences]\label{lem:colored-image-order}
Let \(X\) be a set of isotropic elements of \(G_B\) in a single
overall color, and let
\[
 H=\langle x-y:x,y\in X\rangle\subset G_B.
\]
For the restricted pairing map
\[
 \theta_H:H\longrightarrow\Hom(H,\mathbb Q/\Z),\qquad
 \theta_H(u)=\lambda_B(u,\cdot)|_H,
\]
one has
\[
 |\im\theta_H|\le N^{NJ_B}.
\]
\end{lemma}

\begin{proof}
Fix a prime and abbreviate \(h=b_{p,3}\). In the orthogonal
decomposition \(G_{B,p}=U\oplus T\) used in
Lemma~\ref{lem:primary-coloring}, let
\(d=\#\{i:b_{p,i}>h\}\le2\). Project the \(p\)-primary part of \(H\)
to subgroups \(H_U\subset U\) and \(H_T\subset T\).
Every pairing between projections of elements of \(X\) is
annihilated by \(p^h\) on \(U\). Bilinearity gives the same
conclusion for pairings of arbitrary elements of \(H_U\), since
these are generated by projections of differences.

A subgroup of a finite abelian \(p\)-group with \(d\) cyclic factors
has at most \(d\) generators. One way to see this is that its
minimal number of generators equals the dimension of its subgroup
annihilated by \(p\), which is at most
\(\dim_{\mathbb F_p}U[p]=d\). The image of the restricted pairing
map on \(H_U\) is a quotient of \(H_U\), and it is annihilated by
\(p^h\). Its order is therefore at most \(p^{hd}\). The image
order of the restricted pairing on \(H_T\) is at most
\(|H_T|\le|T|\).

Pull these two pairings back to the \(p\)-primary part of \(H\).
The image of the sum of their maps is contained in the sum of
their images, whose order is at most the product of the two
orders. Consequently the image order at this prime is at most
\[
 |T|p^{hd}
 =p^{\sum_i\min(b_{p,i},h)}
 =p^{\sum_{e\le b_{p,3}}k_B(p,e)}.
\]
The range \(e\le b_{p,3}\) is exactly the range in which
\(k_B(p,e)\ge3\). Distinct primary pairings are orthogonal, so
multiplying the last bound over primes gives
\[
 |\im\theta_H|
 \le\prod_p p^{\sum_{e:\,k_B(p,e)\ge3}k_B(p,e)}
 =N^{NJ_B},
\]
as required.
\end{proof}

\subsection{The graph lattice and the spectral comparison}

We first record the determinant consequence of a pairing image
bound. All lattice determinants in this section are Euclidean
covolumes in their real spans.

\begin{lemma}[Pairing index and determinant]\label{lem:pairing-determinant}
Let \(L\) be a lattice, let \(V=\Span L\), and let \(q\) be a
symmetric bilinear form on \(V\) such that \(q(L,L)\subset\mathbb Q\).
Write
\[
 K=L\cap\operatorname{rad}q,\qquad
 S=V\cap(\operatorname{rad}q)^\perp,\qquad
 \bar L=\operatorname{proj}_S L.
\]
Then \(K\) spans the radical and is primitive in \(L\), and
\(\bar L\) is a lattice in \(S\). If \(t\) is the order of the
image of
\[
 L\longrightarrow\Hom(L,\mathbb Q/\Z),\qquad
 x\longmapsto q(x,\cdot)\bmod\Z,
\]
then
\begin{equation}\label{eq:pairing-determinant}
 \left|\det_{\mathrm{orth}}(q|_S)\right|
 \ge \frac{1}{t(\det\bar L)^2}
 =\frac{(\det K)^2}{t(\det L)^2}.
\end{equation}
Here the determinant on the left is taken in an orthonormal basis
of \(S\), and all rank-zero determinants are one.
\end{lemma}

\begin{proof}
The matrix of \(q\) in an \(L\)-basis is rational. Its kernel is
therefore rational in that basis, so \(K\) spans the radical.
It is primitive because it is the intersection of \(L\) with a
real subspace. Extending a basis of \(K\) to a basis of \(L\)
shows both that \(\bar L\) is a lattice and that
\[
 \det L=\det K\,\det\bar L.
\]
The form vanishes whenever one argument is in the radical, so it
descends to the projected lattice \(\bar L\). Its pairing map
there has the same image order \(t\): the quotient \(L\to\bar L\)
is surjective, and pullback injects
\(\Hom(\bar L,\mathbb Q/\Z)\) into \(\Hom(L,\mathbb Q/\Z)\).
The image is finite because the form has rational entries in a
lattice basis.

Let \(M\) be its matrix in a basis of \(\bar L\), of rank
\(s=\dim S\). It is a nonsingular rational matrix. The kernel of
the pairing modulo \(\Z\) in \(\Z^s\) has index \(t\); choose its
integer basis matrix \(A\), so \(|\det A|=t\).
The matrix \(A^{\mathsf T}M\) is integral. It has nonzero
determinant, and hence
\[
 1\le |\det(A^{\mathsf T}M)|=t|\det M|.
\]
This uses the pairing index on one side of the matrix. If \(P\)
is the basis matrix of \(\bar L\) in orthonormal coordinates on
\(S\), then \(M=P^{\mathsf T}M_{\mathrm{orth}}P\) and
\(|\det P|=\det\bar L\). Thus
\[
 |\det M_{\mathrm{orth}}|
 =\frac{|\det M|}{(\det\bar L)^2}
 \ge\frac{1}{t(\det\bar L)^2},
\]
which proves Equation~\eqref{eq:pairing-determinant}. The rank-zero
case has \(t=1\) and follows from the stated conventions.
\end{proof}

\begin{proof}[Proof of Proposition~\ref{prop:admissible-potential}]
Fix \(A>0\), and choose a fixed \(P>A+2\). Intersect the typical
event of Corollary~\ref{cor:inverse-rank} with exponent \(D=2\) in the
subspace height bound and precision \(P\), and that of
Proposition~\ref{prop:typical-spectrum} with precision \(P\).
This is a typical event uniformly for \(m\in\mathcal I_N\).
On it the rank conclusion holds simultaneously for every rational
subspace \(W\subset\R^m\) with \(H(W)\le\exp(N\ell^2)\); the
spectral conclusion holds for \(B\) itself. We will verify that
the subspace selected below meets this height bound, so the
selection may depend on \(B\) and on all its admissible columns.

Fix a nonsingular \(B\) in this event. There is nothing to prove
if \(\mathcal S(B)\) is empty. Otherwise choose a largest overall
color among its elements, with the color of a column meaning the
color of its class in \(G_B\). Lemma~\ref{lem:robust-extraction},
applied to this color with parameter \(\rho=\ell^{-P}\), gives a
nonempty subset \(T\) robust on its affine span and satisfying
\begin{equation}\label{eq:discriminant-size-loss}
 \log_2|\mathcal S(B)|
 \le \log_2 C_{\mathrm{col}}(B)+N\ell^{-P}+\log_2|T|.
\end{equation}
Put \(W=\Span(T-T)\) and \(r=\dim W\). This is a rational
subspace, since it is spanned by integer differences. Projection
onto a suitable set of \(r\) coordinates is injective on \(W\),
and hence on the affine span of \(T\). Therefore
\begin{equation}\label{eq:discriminant-cube-dimension}
 |T|\le2^r.
\end{equation}

The graph below uses admissibility to keep differences short and retains
integral first coordinates for the determinant bound on the radical lattice.
For an eigenvalue \(\lambda\) of \(B\), its induced metric changes the
corresponding eigenvalue \(1/\lambda\) of the inverse form to
\(\lambda/(1+\lambda^2)\), which is bounded even when \(\lambda\) is small.
Consider the rational injective graph map
\[
 \gamma:\R^m\longrightarrow\R^{2m},\qquad
 \gamma(z)=(z,B^{-1}z),
\]
and the rank-\(r\) lattice
\[
 L=\langle\gamma(z)-\gamma(t):z,t\in T\rangle_{\Z}
 \quad\text{in }\gamma(W).
\]
Each graph difference has squared norm at most \(m+4m\le5N\):
the first coordinates are differences of bits, and each inverse
image has infinity norm at most one. The restriction of
\(\gamma\) to the affine span of \(T\) is a rational affine
injection, and it preserves the codimensions occurring in
robustness. Lemma~\ref{lem:robust-height} therefore applies to
these graph points. It also applies to the original cube points.
For a constant \(C_0\) depending only on \(P\), it gives
\begin{equation}\label{eq:discriminant-lattice-height}
 \det L\le\exp(C_0N\ell),\qquad
 H(W)\le\exp(C_0N\ell).
\end{equation}
For all sufficiently large \(N\), the second bound in
Equation~\eqref{eq:discriminant-lattice-height} is at most
\(\exp(N\ell^2)\). Thus this adaptively chosen \(W\) lies in the
simultaneous scope of Corollary~\ref{cor:inverse-rank}.
Lemma~\ref{lem:robust-height} is deterministic for every such
robust set and graph map, and the spectral event is a statement
about all eigenvalues of \(B\); neither requires \(W\) to have
been fixed in advance.

On the whole graph \(\Gamma=\gamma(\R^m)\), define
\[
 Q(\gamma(x),\gamma(y))=x^{\mathsf T}B^{-1}y.
\]
Let \(s\) be the rank of its restriction to \(\gamma(W)\).
The graph map identifies that restriction with the restriction
of the bilinear form \(B^{-1}\) to \(W\). Hence
Corollary~\ref{cor:inverse-rank} gives
\begin{equation}\label{eq:discriminant-rank}
 s\ge r-N\ell^{-P}.
\end{equation}
The reduction of \(Q\) modulo \(\Z\) on \(L\) is the pullback of
the restricted discriminant pairing on
\[
 H=\langle \bar z-\bar t:z,t\in T\rangle\subset G_B.
\]
Indeed, the map \(L\to H\) sending \(\gamma(u)\) to
\(u+B\Z^m\) is surjective and respects these pairings. Pullback
on the corresponding character groups is injective, so the two
pairing images have the same order. By
Lemma~\ref{lem:colored-image-order}, this order is at most
\(N^{NJ_B}\).

Let \(K=L\cap\operatorname{rad}(Q|_{\gamma(W)})\), and let \(S\)
be the orthogonal complement of this radical in \(\gamma(W)\).
The rationality assertion in Lemma~\ref{lem:pairing-determinant}
applies since \(B^{-1}\) is rational and the first coordinates
of vectors in \(L\) are integral. Consequently \(\rank L=r\),
\(\rank K=r-s\), and \(\dim S=s\). There is also the lower bound
\(\det K\ge1\). In fact the projection to the first \(m\)
coordinates is injective on \(\Gamma\), and it maps \(K\) to an
integer lattice of the same rank. The squared determinant of an
integer lattice is the positive integer determinant of the Gram
matrix of an integer basis, so its determinant is at least one.
Orthogonal projection contracts Euclidean volumes, giving the
same lower bound for \(\det K\). This includes rank zero by
convention.

Lemma~\ref{lem:pairing-determinant}, followed by
Equation~\eqref{eq:discriminant-lattice-height}, now gives
\begin{equation}\label{eq:discriminant-det-lower}
 \left|\det_{\mathrm{orth}}(Q|_S)\right|
 \ge N^{-NJ_B}(\det L)^{-2}
 \ge N^{-NJ_B}\exp(-2C_0N\ell).
\end{equation}

We compare this with the spectrum of \(Q\) on \(\Gamma\), using
the Euclidean metric induced from \(\R^{2m}\). If \(v_i\) are
orthonormal eigenvectors of \(B\), with nonzero eigenvalues
\(\lambda_i\), then their graph vectors are orthogonal and
\[
 \|\gamma(v_i)\|^2=1+\lambda_i^{-2}.
\]
It follows that the self-adjoint operator representing \(Q\)
on \(\Gamma\) has eigenvalues
\[
 \mu_i=\frac{\lambda_i}{1+\lambda_i^2}.
\]
All have absolute value at most one. Let \(C_P\) be the fixed
constant supplied by Proposition~\ref{prop:typical-spectrum}.
Apart from \(t_B\le N\ell^{-P}\) eigenvalues, that proposition
gives \(|\lambda_i|>\sqrt N\,\ell^{-C_P}\). For these indices,
\[
 |\mu_i|\le|\lambda_i|^{-1}\le a_N,\qquad
 a_N=\frac{\ell^{C_P}}{\sqrt N}<1
\]
for sufficiently large \(N\).

The determinant of a compression to any \(s\)-dimensional
subspace is bounded in absolute value by the product of the
largest \(s\) singular values of the whole operator. Indeed, if
\(w\) is the unit exterior product of an orthonormal basis of
that subspace and \(T_Q\) is the operator of \(Q\), the
determinant is
\(\langle w,(\bigwedge^s T_Q)w\rangle\), whose absolute value
is at most \(\|\bigwedge^s T_Q\|\). Apply this to the subspace
\(S\). Using the spectral bounds above, including the rank-zero
convention, gives
\begin{equation}\label{eq:discriminant-det-upper}
 \left|\det_{\mathrm{orth}}(Q|_S)\right|
 \le a_N^{(s-t_B)_+}\le a_N^{s-t_B}.
\end{equation}
The last inequality is also valid for \(s<t_B\), since \(a_N<1\).
Comparing Equations~\eqref{eq:discriminant-det-lower}
and~\eqref{eq:discriminant-det-upper} and taking logarithms yields
\[
 (s-t_B)\left(\frac12\log N-C_P\log\ell\right)
 \le NJ_B\log N+2C_0N\ell.
\]
The coefficient on the left is positive for large \(N\).
Together with Equation~\eqref{eq:discriminant-rank} and
\(t_B\le N\ell^{-P}\), this implies
\[
 \frac rN
 \le
 \frac{J_B+2C_0\ell/\log N}
      {1/2-C_P\log\ell/\log N}
       +2\ell^{-P}
 =2J_B+O(\ell^{-P}).
\]
For the last equality, use \(J_B\le1/2\), the fixedness of
\(C_0,C_P,P\), and
\((\ell+\log\ell)/\log N=o(\ell^{-P})\).

Finally combine Equations~\eqref{eq:discriminant-size-loss},
\eqref{eq:discriminant-cube-dimension}, and
\eqref{eq:discriminant-color-cost}. They give
\[
 \begin{split}
 \frac{\log_2|\mathcal S(B)|}{N}
 &\le
 \beta\sum_{\substack{p,e\\k_B(p,e)=2}}\nu_p
       +2J_B+O(\ell^{-P})+O(N^{-1/5}\log N)\\
 &=\Psi(B)+O(\ell^{-P})+O(N^{-1/5}\log N).
 \end{split}
\]
The equality follows from Equations~\eqref{eq:potential}
and~\eqref{eq:discriminant-high-weight}. Because \(P>A+2\) is
fixed, the sum of the two error terms is at most \(\ell^{-A}\) for all
sufficiently large \(N\). This proves the claimed estimate on
the chosen typical event and completes the proof.
\end{proof}
